\documentclass[11pt, a4paper]{article}

% Paquets de base
\usepackage[utf8]{inputenc}
\usepackage[T1]{fontenc}
\usepackage[french]{babel}
\usepackage[margin=2.5cm]{geometry}
\usepackage{hyperref}
\usepackage{enumitem}

% Informations du document
\title{\textbf{Compte Rendu de Réunion de Lancement - TER}}
\author{ABBES Samy \& AREZKI Celia}
\date{Réunion du 4 Mai 2026}

\begin{document}

\maketitle

\section*{Informations Générales}
\begin{itemize}
    \item \textbf{Date de la réunion :} 4 Mai 2026 à 14h
    \item \textbf{Membre du groupe présents :} ABBES Samy, AREZKI Celia
    \item \textbf{Encadrants présents :} Dupé François-Xavier, LABOUREL Arnaud
    \item \textbf{Sujet du TER :} Benchmarking et analyse comparative d'algorithmes d'agrégation robustes en apprentissage fédéré face aux adversaires byzantins avec ByzFL
\end{itemize}

\vspace{0.5cm}

\section{Ordre du jour et objectifs de la réunion}
Cette première réunion avait pour but de définir le cadre de travail, les outils à mettre en place pour le suivi du projet, ainsi que les premiers objectifs techniques concernant la prise en main de la bibliothèque d'apprentissage fédéré robuste (\textit{ByzFL}).

\section{Directives et tâches assignées (Court terme)}
Lors de cette séance, les encadrants ont défini les tâches suivantes à réaliser pour la prochaine réunion :
\begin{itemize}
    \item \textbf{Environnement de travail :} Création d'un dépôt Git pour le projet. Les deux professeurs encadrants doivent y être ajoutés avec le rôle de \textit{Reporter}.
    \item \textbf{Prise en main de ByzFL :} Installation de la bibliothèque \textit{ByzFL} et étude approfondie de son fonctionnement (compréhension globale des mécanismes d'attaques byzantines et des méthodes d'agrégation robuste proposées par l'outil) en vue d'une présentation lors de notre prochaine entrevue.
    \item \textbf{Expérience préliminaire :} Mise en place d'une petite expérience utilisant \textit{ByzFL}, qui devra être exécutée et démontrée lors de la prochaine réunion.
    \item \textbf{Documentation initiale :} Rédaction du présent compte rendu et préparation d'un premier rapport de synthèse résumant le travail d'installation et l'expérience réalisée.
\end{itemize}

\section{Règles de fonctionnement et livrables (Long terme)}
Pour la suite du Travail d'Étude et de Recherche, les règles de suivi suivantes ont été actées :
\begin{itemize}
    \item \textbf{Comptes rendus :} Un compte rendu de réunion doit être rédigé et transmis au plus tard \textbf{2 jours} après chaque réunion.
    \item \textbf{Suivi hebdomadaire :} Un rapport hebdomadaire résumant le travail effectué durant la semaine doit être fourni.
    \item \textbf{Code et reproductibilité :} Le code des expériences effectuées devra être systématiquement fourni, accompagné d'une description précise de l'environnement Python nécessaire (\texttt{requirements}) à leur mise en œuvre, afin de faciliter et d'assurer la reproductibilité des expériences menées.
\end{itemize} 

\section{Prochaines étapes}
\begin{itemize}[label=$\square$]
    \item Créer le dépôt Git et inviter les encadrants.
    \item Installer \textit{ByzFL} et préparer la présentation de l'outil.
    \item Coder l'expérience de démonstration.
\end{itemize}

\end{document}
